\hypertarget{conclusion}{%
\section{Conclusion}\label{conclusion}}

\hypertarget{summary-of-work}{%
\subsection{Summary of work}\label{summary-of-work}}

This project set out to answer a specific question. What changes when a
cascade router is given live infrastructure state and a learned, closed
decision loop, instead of a confidence threshold fixed at design time?
The artefact built to answer it is a context-aware agentic gateway with
seven integrated components. These are the Edge Context Protocol, the
dynamic Medallion pipeline, an MCP-pattern skill interface, a PPO agent
trained offline in a Gymnasium simulation, a Cascade Controller wiring
the request path together, a bidirectional forward/backward loop, and a
Policy Orchestration Layer governing multiple clients' policies. It was
evaluated against five core baselines and eight literature baselines
(Section~\ref{literature-survey}), across three traffic scenarios, plus the full seven-run
ablation (Section~\ref{evaluation-and-discussion}).

The headline result is the bursty-traffic comparison (Section~\ref{baselines-and-scenarios}).
Every static-threshold or always-cloud baseline that tries to preserve
accuracy under a traffic burst pays for it with SLA violations in the
59.5--60.5\% range. The proposed system, in contrast, holds zero SLA
violations, at 0.5\% of the escalation rate and 14\% of the energy cost
of \textbf{always\_cloud}, and within two percentage points of its
accuracy.

Under steady traffic the accuracy question is closer. A simple static
threshold reaches 0.805 against the proposed system's 0.785, so it represents
two accuracy points when nothing is stressed. But it
buys those two points at 3.5 times the energy, and by escalating every
request instead of just 1\% of them. The distinctive claim is still the
stressed condition.

\hypertarget{revisiting-the-objectives}{%
\subsection{Revisiting the objectives}\label{revisiting-the-objectives}}

The seven objectives set out in Section~\ref{aims-and-objectives} were addressed as follows.
The literature survey (Objective 1, Section~\ref{literature-survey}) positions the artefact
against RouteLLM, the NSGA-II router, Splitwise, GreenServ and PROTEUS,
as planned. The escalation decision was formalised as a Markov Decision
Process over a 13-slot state, three actions, and a four-term reward
(Objective 2). The Bronze/Silver/Gold Medallion pipeline, the Edge
Context Protocol, and the MCP-pattern skill interface were all
implemented. All three are exercised by the live gateway as well as
the offline simulation (Objectives 3, 4 and 6).

The PPO agent was trained offline with domain randomisation, and
deployed as a frozen policy reachable over HTTP (Objective 5). It has
been live-tested twice, once as a plain-Docker multi-region GCP
deployment and once as a genuine 3-region Kubernetes (GKE) deployment
running the current policy. As Section~\ref{design-development-and-implementation}'s Deployment subsection states
plainly, neither of these is a literal sidecar.

Objective 7 went beyond its original plan. The system was evaluated
against the five core baselines across three traffic scenarios, plus
the full seven-run ablation (Section~\ref{evaluation-strategy}). Eight literature baselines
and a governance-layer ablation (Run 7) were added once the Policy
Orchestration Layer entered scope, in Week 7--8. That addition traces
back to specific supervisor feedback, instead of being presented here
as if it had always been the plan.

Five of the six contributions claimed in Section~\ref{contributions} are backed by
specific measurements. The sixth is not, and that is reported here.

Ablation Run 3 (Section~\ref{ablation}) adds \textbf{calibration\_delta} to a
confidence-only threshold, and that alone makes almost no difference.
Run 1, the full PPO policy, uses the same calibration state and does
much better. Together these point to the Edge Context Protocol's value
coming from feeding a learned policy instead of from adding one more
input to a fixed threshold. Run 3 does not prove this by itself; it
rules out the simplest alternative explanation.

The bidirectional loop is the contribution the final evaluation does
not support. Under the pre-retrain policy, Ablation Run 4 showed that
disabling the loop cost 5.5 accuracy points under degraded network,
exactly the condition it was designed for. Under the retrained policy,
that gap closed. Run 4 now matches \textbf{run1\_full} to within half
an accuracy point in every scenario (Section~\ref{ablation}). The loop's measured
value was real, but it was largest when the policy was least well
adapted to its own inputs. A policy retrained under a fitted classifier
recovers roughly 91\% of what the loop had been supplying (the 5.5-point
gap closes to 0.5 points). So the
architectural case for closing the loop still stands on Section~\ref{literature-survey}'s
survey gap, but the empirical case does not survive the retraining process, and it
is not claimed here.

The policy's learned response to per-request SLA-remaining and
predicted-RTT signals is the mechanism behind the bursty-traffic result
above, shaped by the reward function's SLA term. It is a learned
preference rather than a hard pre-dispatch gate, as mentioned in both
Section~\ref{design-development-and-implementation} and Section~\ref{evaluation-strategy}.

Ablation Run 7 (Section~\ref{ablation-run-7}) shows the Policy Orchestration Layer's
per-client onboarding decisions against a single shared policy. The
LLM-backed governance review itself was not re-run against the
retrained policy; its one escalation of an under-provisioned decision,
one the deterministic rule alone would have missed, comes from the
earlier archived run.

The dynamic Medallion pipeline (Contribution 2) is measured
independently, by the three-way static/conditional/agent-driven Silver
comparison, Run 5 (Section~\ref{evaluation-strategy}). The differences there are real but
modest and scenario-dependent, and no one Silver strategy dominates
throughout, which is itself the finding.

The MCP skill interface (Contribution 4) is benchmarked independently
against real HuggingFace inference latency in the live request path.
Perception cost (Bronze to Silver to Gold to
\textbf{AgenticOrchestrator.decide()}, including the skill interface)
averages 0.57ms, with a p95 of 0.97ms, against a 23.9ms mean real
DistilBERT forward pass. That is about 2.4\% of edge inference latency,
confirming it sits well inside the per-request perception budget
claimed in Section~\ref{contributions}.

\hypertarget{limitations-and-future-work}{%
\subsection{Limitations and future work}\label{limitations-and-future-work}}

Section~\ref{limitations} already mentions six specific limits in
detail (trained \textbf{MiscalibrationClassifier}'s own accuracy
being modest, residual degraded-network over-escalation, the
meta-policy never converging, synthetic queue/RTT/energy signals,
concurrent LLM diagnosis calls happening sequentially, and per-request
SLA check being a learned signal, in place of a hard pre-dispatch
gate). They are not written out in their entirety here. The next steps
at a project-wide level are as follows:

\begin{enumerate}
\item The MiscalibrationClassifier is now trained on real,
non-circular data, in both domains to which it ever was deployed
(Section~\ref{design-development-and-implementation}). The
simulation-domain fit is, by default, the basis of every report
evaluation number, since that point. Though its cross-validated
accuracy is still quite moderate, in particular on the STRESSED class
(33\% precision), the normal next action is not fitting a classifier
for the first time, but improving the fit that already exists. The two
biggest means of action are gathering a bigger simulation-domain
sample and using a better set of features than the six-dimensional
stress-plus-error rate vector that it currently views.

\item The synthetic queue-depth
and RTT signals in the Gymnasium simulation will be replaced by those
extracted from the live multi-region Docker deployment to close the
one sim-to-real gap which the project has not yet been able to
directly validate. The inference latency is already measured on real
hardware, so only this part of the signal remains.

\item The residual 1.5\% over-escalation under sustained network
degradation (Section~\ref{limitations}) is now small enough that the reward function
is no longer the obvious place to attack it. Fitting the classifier
already cut this from 21.5\%, without touching the reward at all
(Section~\ref{misc-classifier-refit}). So the remaining question is whether a
scenario-conditioned reward term would close the last 1.5\%, or simply
trade it for a regression elsewhere. That is a hypothesis worth
testing, as opposed to a diagnosed defect waiting on a known fix, and it should
be evaluated the same way the reward-gradient fix was.

\item LiteLLM's and Portkey's routing strategies (Section~\ref{literature-survey}) are
compared here only qualitatively, against their own documentation.
Turning that into an empirical row in Section~\ref{evaluation-and-discussion} first requires
deciding, fairly to both sides, how a provider-routing gateway's rule
maps onto an edge/cloud decision. Section~\ref{evaluation-strategy} states this plainly,
without presenting an invented mapping as already settled. Carrying this out is the most direct way of reinforcing this paper's
main line of argument - that the system is better, not only compared
to the real production gateways but also against the paper baselines
that were tested.

\item \textbf{sla\_violation\_predicted}
(Section~\ref{design-development-and-implementation}, Silver) is a
signal computed by the system which no decision path accesses at the
moment. Plugging it into \textbf{AgenticOrchestrator.decide()} as an
unconditional pre-dispatch check would transform the project's
current soft SLA preference into a hard per-request guarantee which
was originally proposed.

\item The Gold vector's 13 slots are hard-coded in the code. To add
a new signal, one has to modify \textbf{DynamicMetricRegistry.snapshot()}'s
match-case, add a new field to \textbf{SilverFeatureVector}, and
manually extend \textbf{GoldStateVector.from\_silver()}
(Section~\ref{design-development-and-implementation}). A cheap
partial fixing way is pre-allocating some reserved, zero-masked slots
beyond the 13 used at present, using the same mask already applied to
\textbf{calibration\_delta} and \textbf{error\_rate}. Afterwards, a
new metric could be linked to one of those slots by a name-to-slot
look-up table without touching the source code. With this approach,
it will support all the newly introduced metrics until one day all
the reserved slots have been filled up. Beyond that stage, the
restriction will shift over to the fixed input layer of
\textbf{PPONetwork} (\textbf{agent.py}), which can not be extended
without modifying and resizing its first weight matrix and
restarting training from a clean slate. So this direction
buys a fixed number of free additions. It does not remove the ceiling.
\end{enumerate}

\hypertarget{closing-remarks}{%
\subsection{Closing remarks}\label{closing-remarks}}

The research question posed in Section~\ref{problem-statement} asked what a
context-aware, closed-loop orchestration system buys over a static,
passive threshold. The bursty-traffic figures set out above give a
specific answer to that question, and they hold under conditions no
confidence-only baseline can even observe. The gain is not uniform
across every condition, and this report does not claim that it is. But
it is largest exactly when infrastructure is under the kind of stress
that confidence-only systems have no way to see happening at all.